\documentclass[11pt,a4paper]{article}

\usepackage[american]{babel}
\usepackage[T1]{fontenc}
\usepackage[utf8]{inputenc}
\usepackage{lmodern}
\usepackage{microtype}
\usepackage[a4paper,margin=24mm,headheight=14pt]{geometry}
\usepackage{graphicx}
\usepackage{booktabs}
\usepackage{tabularx}
\usepackage{longtable}
\usepackage{enumitem}
\usepackage{xcolor}
\usepackage{fancyhdr}
\usepackage{hyperref}
\usepackage{libertine}

\definecolor{avacblue}{HTML}{1F5C99}
\definecolor{avaclight}{HTML}{EAF3FA}
\definecolor{avacgray}{HTML}{4E5963}
\hypersetup{
  colorlinks=true,
  linkcolor=avacblue,
  urlcolor=avacblue,
  pdftitle={AVAC4QGIS User Interface Guide},
  pdfauthor={LHE--EPFL},
  pdfsubject={AVAC4QGIS 0.6.1 user interface reference}
}
\setlist[itemize]{leftmargin=*,itemsep=3pt,topsep=4pt}
\setlist[enumerate]{leftmargin=*,itemsep=5pt,topsep=4pt}
\setlength{\parindent}{0pt}
\setlength{\parskip}{5pt}
\renewcommand{\arraystretch}{1.15}

\pagestyle{fancy}
\fancyhf{}
\fancyhead[L]{\small AVAC4QGIS User Interface Guide}
\fancyhead[R]{\small Version 0.6.1}
\fancyfoot[C]{\thepage}

\newcommand{\control}[1]{\textbf{#1}}
\newcommand{\pathcode}[1]{\texttt{#1}}
\newcommand{\uiscreenshot}[2]{%
  \begin{figure}[p]
    \centering
    \includegraphics[width=\textwidth,height=0.78\textheight,keepaspectratio]{images/#1}
    \caption{#2}
  \end{figure}
}
\newenvironment{notebox}{%
  \begin{quote}\small\color{avacgray}\noindent
}{%
  \end{quote}
}

\title{\textbf{AVAC4QGIS User Interface Guide}\\[6pt]
       \large Interface reference for AVAC4QGIS 0.6.1}
\author{LHE--EPFL, Switzerland}
\date{September 2, 2026}

\begin{document}
\maketitle
\thispagestyle{empty}

\begin{notebox}
This guide describes the controls that are visible in the packaged QGIS
plugin. It was checked against the current implementation in
\pathcode{avac\_qgis/gui/dock.py} and the preparation, coupling, and result
modules used by those controls. Hidden compatibility fields and developer-only
code paths are not presented as user features.
\end{notebox}

\tableofcontents
\clearpage

\section{General description of the interface}

AVAC4QGIS is a dockable QGIS panel for preparing, running, and inspecting
depth-averaged avalanche simulations. The standard workflow has three tabs:
\control{AVAC Parameters}, \control{AVAC Run}, and \control{Results}. Inputs,
scientific parameters, execution, visualization, profiles, gauges, and PNG
exports remain inside QGIS, while packaged native solvers run in isolated
case directories.

The optional \control{Enable Lake-Wave Extension} check box inserts
\control{WAVE Parameters} and \control{WAVE Run}. WAVE consumes a completed
AVAC run as a read-only source and uses the same final Results tab. Enabling
WAVE does not change an existing AVAC run.

The parameter and result pages use accordion controls. Click an accordion
header to display its contents. Each major operation reports its state through
status text and a progress bar. Long execution logs remain hidden until their
corresponding log button is enabled.

\uiscreenshot{ui_overview.png}{The complete dock with the persistent controls above the tabs and the normal AVAC workflow.}

\section{General simulation workflow}

\subsection{AVAC-only simulation}

\begin{enumerate}
  \item Choose a writable \control{AVAC Working Directory} above the tabs.
  \item In \control{AVAC Parameters}, select the terrain DEM and release
        polygons, then load or set the release, rheology, and numerical values.
  \item In \control{AVAC Run}, use \control{Check Environment} and
        \control{Validate Inputs}. Use the previews if the initial state or
        altitude-dependent rheology must be inspected.
  \item Select \control{Prepare AVAC Run}. Preparation creates a new,
        self-contained run below \pathcode{runs/}; it does not execute the
        solver.
  \item Select \control{Run}. A completed run becomes available in
        \control{Results}.
  \item In \control{Results}, select the run and load summary or temporal
        products. Profiles, gauges, and exports load a missing temporal product
        automatically before continuing.
\end{enumerate}

\subsection{AVAC followed by WAVE}

\begin{enumerate}
  \item Complete AVAC first and enable the Lake-Wave Extension.
  \item In \control{WAVE Parameters}, choose the completed AVAC run, terrain
        DEM, lake water-body polygon, water-surface elevation, WAVE cell size,
        and model settings. The lake polygon may instead be created from the
        DEM, water elevation, and a map click before AVAC has been run.
  \item Use \control{Preview Water Level}. The exact solver-grid lake state is
        cached and reused by preparation while its inputs remain unchanged.
  \item In \control{WAVE Run}, validate and prepare the WAVE scenario. WAVE
        inherits the completed AVAC calculation bounds, duration, output
        count, and temporal origin. Preparation creates the internal shoreline
        source without modifying AVAC output.
  \item Run WAVE and return to the shared Results tab. Select both the completed
        AVAC run and its completed WAVE scenario when comparing snow and water
        products.
\end{enumerate}

\begin{notebox}
WAVE and AVAC both solve on rectangular GeoClaw grids, but the snow-to-water
transfer is not imposed at the outer rectangle. It is applied internally on
the grid-aligned shoreline of the initially wet lake. Only AVAC flux directed
into the wet lake contributes to the WAVE source.
\end{notebox}

\section{Function of each interface element}

This section follows the current visual order of the interface. It describes
what each visible element reads, changes, or creates.

\subsection{Controls above the tabs}

\begin{itemize}
  \item \control{AVAC Working Directory} is the required writable workspace.
        Type a path or use the browse button. It is user data storage, not the
        plugin installation or the managed runtime. QGIS remembers the path.
  \item \control{Enable Lake-Wave Extension} shows or hides WAVE Parameters
        and WAVE Run and adds WAVE variables and Lake Volume to Results. It
        does not alter AVAC files.
  \item \control{Load Case} restores a complete AVAC4QGIS case YAML: workspace,
        persistent input-layer sources, visible AVAC values, WAVE enabled
        state, and WAVE setup when present.
  \item \control{Save Case} writes the corresponding complete case YAML. It
        stores references to persistent layer sources; temporary memory layers
        must be saved to disk before the case can be portable.
  \item \control{Help} opens this PDF. The Plugins $>$ AVAC4QGIS menu also
        contains Help. Its online-documentation and website entries are present
        but disabled until those destinations exist.
\end{itemize}

\subsection{AVAC Parameters tab}

This tab contains four accordion pages and the AVAC-only configuration
buttons. Changing a visible input or parameter invalidates any prepared run;
prepare again before executing it.

\subsubsection{AVAC Inputs}

\begin{itemize}
  \item \control{DEM raster layer} selects the fixed basal topography. The
        first raster band must have a valid CRS, finite usable cells, and square
        pixels. The selected DEM extent defines the rectangular AVAC
        calculation domain. The terrain resolution must equal or evenly
        subdivide the chosen computational cell size, and the selected DEM
        dimensions must be divisible by the corresponding number of DEM cells
        per computational cell. AVAC4QGIS rejects an incompatible extent
        rather than silently cropping release cells or extending the solver
        domain. Internally it writes one replicated edge-sample halo for
        GeoClaw terrain interpolation only; this does not add release or qinit
        cells. If the optional fine-DEM backend is enabled programmatically,
        its local terrain file uses a closed node envelope at the selected fine
        DEM edges, so it cannot replace the base terrain outside that footprint.
  \item \control{Release polygon layer} selects polygon or multipolygon
        features defining the initially mobile release area. Geometries are
        transformed to the DEM CRS during preparation when needed. At least one
        DEM grid point must fall inside the combined polygons.
\end{itemize}

Preparation copies the selected sources into the run's \pathcode{inputs/}
tree for traceability. The optional refinement-DEM backend remains in the
source for compatibility but has no visible selector and is not part of the
normal workflow.

\uiscreenshot{ui_avac_parameters_inputs.png}{AVAC Parameters with the AVAC Inputs accordion open.}

\subsubsection{Release / initial conditions}

The DEM remains the fixed basal surface. These controls initialize only the
mobile release depth. There is no additional stationary winter snowpack:
outside the release polygons, initial mobile depth is zero and the initial
snow-surface product equals the DEM.

\begin{itemize}
  \item \control{Release depth} $d_0$ (m) is the base mobile thickness. With
        both corrections disabled, every release-grid point receives $d_0$.
  \item \control{Critical slope} $\theta_{cr}$ (degrees) and
        \control{Slope correction} $\nu$ define the optional local-slope
        factor. The implementation evaluates it only where DEM slope exceeds
        $25^\circ$ and the denominator is numerically safe.
  \item \control{Reference elevation} $z_{ref}$ (m) and
        \control{Hypsometric gradient} define the optional elevation change
        $(z-z_{ref})\,g_h/100$. A positive gradient increases depth above
        $z_{ref}$ and decreases it below $z_{ref}$.
  \item \control{Apply slope correction} and \control{Apply elevation
        correction} select which corrections are applied. If both are active,
        the elevation-corrected depth is multiplied by the slope factor.
  \item \control{Return period [years]} is positive integer scenario metadata.
        It labels the case and results but does not change AVAC equations or
        the release-depth calculation.
\end{itemize}

The initial-depth preview should be inspected when corrections are active.
AVAC4QGIS rejects non-finite release controls, a release polygon over a
non-finite DEM cell, and any correction that produces a negative or non-finite
candidate depth. It does not silently replace those values with zero: revise
the release depth, reference elevation, hypsometric gradient, or slope
correction before preparing the run.

\uiscreenshot{ui_avac_parameters_release.png}{AVAC Parameters with Release / initial conditions open.}

\subsubsection{Rheology}

\begin{itemize}
  \item \control{Model} selects \control{Coulomb}, \control{Voellmy}, or
        \control{cohesive\_Voellmy}. Coulomb uses basal dry-friction
        resistance. Voellmy adds speed-dependent resistance. Cohesive Voellmy
        adds cohesive strength and its static-yield contribution.
  \item \control{Density} $\rho$ (kg/m$^3$) is bulk avalanche density. It is
        used by the cohesive term and by the displayed kinetic-pressure result
        $p_k=\rho v^2/2$, reported in kPa.
  \item \control{$\mu$} is the dimensionless Coulomb-friction coefficient.
        Increasing it strengthens the speed-independent basal resistance.
  \item \control{$\xi$} (m/s$^2$) controls the Voellmy speed-dependent term.
        In the implemented term, larger $\xi$ means weaker speed-dependent
        resistance. It is disabled for the Coulomb-only model.
  \item \control{Cohesion} $C$ (Pa) adds material strength in the cohesive
        Voellmy model and is disabled for the other two models.
  \item \control{Altitude zones} replaces uniform $\mu$, $\xi$, and $C$ with
        bed-elevation zones. Enter one comma-separated row per zone as
        \pathcode{mu, xi, cohesion, lower elevation}. The first row must leave
        the fourth field blank. Later lower elevations must be strictly
        increasing; a cell exactly at a break belongs to the higher zone.
  \item \control{Visualize Rheology Zones} writes
        \pathcode{qgis\_previews/avac\_rheology\_zones.tif} and loads it as a
        categorical raster. Its colors and legend show exactly which zone each
        valid DEM cell will use. It does not prepare or alter a run.
\end{itemize}

The AVAC source update integrates the local Coulomb/cohesive and Voellmy
resistance with an explicit zero-speed arrested state. This is why rheology is
material behavior, not merely result styling.

\uiscreenshot{ui_avac_parameters_rheology.png}{AVAC Parameters with the Rheology accordion open.}

\subsubsection{Simulation / numerical}

\begin{itemize}
  \item \control{Simulation duration} is the final elapsed solver time in
        seconds.
  \item \control{AMR refinement levels} is the maximum number of adaptive mesh
        levels. Additional levels can resolve flagged flow regions more finely
        and increase cost. Level 1 is the unrefined base calculation.
  \item \control{Output interval} is the requested elapsed time between
        outputs. Its default is exactly 1 s. When duration or interval is
        edited, the plugin calculates a whole number of intervals and includes
        both $t=0$ and final time in the temporal raster. A loaded configuration
        retains its explicit output count until one of these timing controls is
        changed.
  \item \control{CFL target} is the normal adaptive time-step target and
        \control{CFL maximum} is the rejection limit. Target must not exceed
        maximum.
  \item \control{Computational cell size} (m) is base solver spacing. It must
        equal or be a whole-number multiple of the DEM pixel size, and both
        domain spans must contain a whole number of solver cells.
  \item \control{Limiter} selects none, minmod, superbee, MC, or van Leer for
        second-order wave limiting. The AVAC default is superbee.
\end{itemize}

There is no boundary selector. Every normally prepared AVAC case uses
extrapolating outer boundaries, so material may leave the rectangular domain.
The prepared YAML also standardizes binary 32-bit output, zero solver
verbosity, and Depth as the default temporal variable.

\uiscreenshot{ui_avac_parameters_numerical.png}{AVAC Parameters with Simulation / numerical open.}

\subsubsection{Load Configuration and Save Configuration}

\control{Load Configuration} reads a complete AVAC YAML schema, restores the
visible values, and makes that file the template for preparation. Unknown
advanced keys remain in the complete template rather than being discarded.
\control{Save Configuration} applies the current visible values and fixed run
choices to that complete schema and writes a new YAML file after validation.
These buttons affect AVAC settings only; Load Case and Save Case above the tabs
cover the complete plugin setup.

\subsection{AVAC Run tab}

\begin{itemize}
  \item \control{Prepared Simulation} summarizes the isolated run, return
        period, duration, output interval, cell size, and rheology after
        preparation.
  \item \control{Active run} is read-only and shows the directory that Run will
        execute.
  \item \control{CPU cores} sets \pathcode{OMP\_NUM\_THREADS} from one to the
        maximum logical cores detected on the computer. The default is the
        maximum.
  \item \control{Check Environment} checks the writable workspace, complete
        built-in or user-selected configuration template, and managed AVAC
        runtime. External Make, gfortran, Python, and CLAW variables are not
        required for packaged execution.
  \item \control{Validate Inputs} checks DEM, release geometry, visible
        parameters, CRS, and grid compatibility without creating a run.
  \item \control{Prepare AVAC Run} creates a timestamped directory under
        \pathcode{runs/}, materializes source inputs, computes the terrain and
        initial condition, writes the run configuration, and stages runtime
        data. Its progress bar reports this background task.
  \item \control{Preview Initial Depth} loads the exact mobile $h_0$ field used
        by preparation. \control{Preview Initial Snow Surface} loads
        $z+h_0$. Both reusable rasters are written below
        \pathcode{qgis\_previews/} in the Working Directory.
  \item \control{Run} starts the prepared native solver. \control{Stop}
        requests termination of that AVAC process. A stopped run remains on
        disk but is not marked as completed.
  \item The execution status and progress bar report readiness and frame
        progress. \control{Show execution log} reveals environment,
        preparation, solver, and error messages.
\end{itemize}

\uiscreenshot{ui_avac_run.png}{The current AVAC Run tab.}

\subsection{WAVE Parameters tab}

This tab exists only while the Lake-Wave Extension is enabled. It has three
accordion pages. WAVE is a separate water solver and does not write into the
selected AVAC run.

\subsubsection{AVAC Source}

\control{Completed AVAC run} lists completed runs in the Working Directory.
\control{Refresh Completed AVAC Runs} refreshes the list. The selected run
provides AVAC frames, CRS, rectangular calculation bounds, duration, output
count, and the shared temporal origin.

WAVE has no independent calculation-domain controls. Its rectangular solver
domain is exactly the selected completed AVAC domain. WAVE may use a different,
coarser cell size, but its terrain must cover these bounds and share the AVAC
CRS.

\uiscreenshot{ui_wave_parameters_source.png}{WAVE Parameters with the AVAC Source accordion open.}

\subsubsection{Terrain and Lake Inputs}

\begin{itemize}
  \item \control{Terrain/bathymetry DEM} is the WAVE bed source. It must use
        square pixels, have a valid CRS, and cover the complete AVAC domain.
        It is commonly the same source DEM used by AVAC, but it may include
        bathymetric information not present in the AVAC topography.
  \item \control{Lake water-body polygon} defines the actual reservoir or lake
        footprint. It determines the initially wet mask and the internal
        shoreline used for AVAC-to-WAVE transfer. It is not the rectangular
        solver domain.
  \item \control{Water level} (m) is absolute initial free-surface elevation in
        the DEM's vertical datum, not water depth. A cell inside the lake is wet
        only when its prepared bed is below this elevation. A numeric value of
        zero is therefore an elevation of 0 m in that datum; it does not mean
        ``empty lake.''
  \item \control{Create Lake Polygon from Map Point} arms one QGIS map click.
        From the clicked DEM cell it finds the four-connected cells at or below
        Water level, expands its DEM read window until the contour closes, and
        writes a persistent GeoPackage under \pathcode{derived\_inputs/}. It
        requires the DEM, water level, map point, and writable Working
        Directory, but no completed AVAC run. A right click cancels selection.
  \item \control{Preview Water Level} creates and loads the exact WAVE-grid
        initial depth. It uses the lake polygon and selected WAVE cell size,
        stabilizes the grid-aligned shoreline, stores the GeoTIFF under
        \pathcode{qgis\_previews/}, and caches the prepared lake state.
  \item \control{Wave grid cell size} (m) must equal the terrain resolution or
        be a larger whole-number multiple. Coarsening reduces runtime and
        memory; unsupported finer or fractional resampling is rejected.
\end{itemize}

The lake polygon is rasterized consistently with the WAVE mesh. When the DEM
is finer, a solver cell becomes lake only when at least half of its fine-grid
coverage is inside the polygon. A narrow shoreline guard is adjusted to the
waterline plus dry tolerance to prevent artificial waves at $t=0$. Exterior
cells are not permanently dry: resolved runup may wet them later.

\uiscreenshot{ui_wave_parameters_lake.png}{WAVE Parameters with Terrain and Lake Inputs open. The point-derived lake polygon does not require a prior AVAC run.}

\subsubsection{WAVE Model Settings}

\begin{itemize}
  \item \control{Snow-to-water damping} multiplies the internally coupled
        AVAC source rates. The permitted interface range is 0 to 0.4; smaller
        values inject less water volume and momentum.
  \item \control{Land Strickler coefficient} and \control{Water Strickler
        coefficient} set roughness on the two bed-elevation sides of the
        water-level break. The water value applies below the initial water
        elevation and the land value at or above it. WAVE passes their
        reciprocals as Manning coefficients, so larger Strickler means weaker
        Manning resistance.
  \item \control{Friction depth limit} is the maximum water depth on which the
        Manning source term is applied.
  \item \control{Dry-depth limit} is the WAVE wet/dry threshold and the small
        elevation offset used by shoreline stabilization.
  \item \control{Wave refinement tolerance} controls GeoClaw wave-based AMR
        flagging. WAVE currently prepares one AMR level in the normal plugin
        workflow; the tolerance remains part of its configuration.
  \item \control{CFL target} and \control{CFL maximum} control adaptive time
        stepping. Target must not exceed maximum.
  \item \control{Limiter} selects none, minmod, superbee, MC, or van Leer. The
        WAVE default is van Leer.
\end{itemize}

\control{Load WAVE Configuration} and \control{Save WAVE Configuration}
restore or write only this WAVE setup: selected AVAC source reference, terrain,
lake polygon, water level, cell size, and model values. They do not change AVAC
Parameters.

\uiscreenshot{ui_wave_parameters_model.png}{WAVE Parameters with WAVE Model Settings open.}

\subsubsection{Internal AVAC-to-WAVE source}

During WAVE preparation, the plugin constructs north/south/east/west faces
around the initially wet lake on the WAVE grid. At each AVAC output time it
samples the three conserved AVAC components $h$, $hu$, and $hv$ at each face.
The inward normal flux is
\[
q_n = hu\,n_x + hv\,n_y.
\]
Only finite samples with positive depth and $q_n>0$ are active. Multiplication
by face length and Snow-to-water damping produces an integrated water-volume
rate $Q$. WAVE receives all three conservative source rates
\[
(Q,\; Q u,\; Q v),
\]
not mass alone. The WAVE solver divides those integrated rates by the area of
the base-grid or AMR cell that receives the point source. This prevents mesh
refinement from changing the total injection.

Preparation integrates the source history and records injected water volume,
both depth-momentum components, and both physical water-momentum components in
\pathcode{CL/summary\_config.yaml} and in the WAVE execution log. If no
inward-moving AVAC sample reaches the wet shoreline, preparation reports that
the lake was not reached and leaves Run unavailable.

\subsection{WAVE Run tab}

\begin{itemize}
  \item \control{Prepared Simulation} identifies the WAVE scenario and its
        AVAC source after preparation.
  \item \control{Check Environment} verifies the writable workspace, managed
        WAVE runtime, solver, dependencies, and selected core count.
  \item \control{Validate Inputs} checks the completed AVAC source, terrain,
        lake polygon, inherited AVAC domain, cell-size alignment, CRS, and
        current water-level preview state without creating a scenario.
  \item \control{Prepare Wave Run} creates a timestamped directory below
        \pathcode{wave\_runs/}, reuses or creates the water-level preview,
        writes terrain and wet/dry inputs, and creates the internal shoreline
        source. WAVE duration and output count are inherited from AVAC.
  \item \control{Run} starts the prepared WAVE solver. \control{Stop} first
        requests a clean termination and force-stops the process after a grace
        period if needed.
  \item \control{CPU cores} sets WAVE's \pathcode{OMP\_NUM\_THREADS} and
        defaults to the detected maximum.
  \item Status text and both progress bars report preparation and output-frame
        progress. \control{Show Wave execution log} reveals the preparation,
        coupling ledger, runtime, and solver messages.
\end{itemize}

\uiscreenshot{ui_wave_run.png}{The current WAVE Run tab.}

\subsection{Results tab}

Results is always the last tab. It is shared by both solvers. Controls select a
product whose label begins with AVAC or WAVE and dispatch the work to the
corresponding completed run. The progress bar remains at the bottom of the tab.

\subsubsection{Results Run}

\begin{itemize}
  \item \control{Completed workspace run} and \control{Refresh Runs} select a
        completed AVAC run from the Working Directory.
  \item \control{Open Run Directory} selects an external completed AVAC run
        directory.
  \item When WAVE is enabled, \control{Completed WAVE scenario}, its Refresh
        Runs button, and \control{Open Wave Scenario Directory} select a
        completed WAVE scenario. AVAC and WAVE summaries remain visible
        together.
\end{itemize}

\uiscreenshot{ui_results_runs.png}{Results Run with the Lake-Wave Extension enabled.}

\subsubsection{Summary Maps}

AVAC summary choices are Maximum Depth, Maximum Velocity, Maximum Pressure,
Arrival Time, Time of Maximum Depth, and Time of Maximum Velocity. Maximum
Pressure is the kinetic-pressure display $\rho v^2/2$ in kPa. Time products
contain elapsed simulation seconds.

When WAVE is enabled, the selector also contains Maximum Surface Rise and
Maximum Surface Drawdown. These are the largest positive rise and largest
positive magnitude of drawdown relative to the initial WAVE free surface.
\control{Load Map} creates or reuses the selected GeoTIFF and adds it to QGIS.

\uiscreenshot{ui_results_summary.png}{The shared Summary Maps accordion. WAVE choices appear in the selector when the extension is enabled.}

\subsubsection{Time Series Maps and PNG export}

\begin{itemize}
  \item \control{Variable} selects AVAC Depth, Velocity, Pressure, or Snow
        Surface Elevation. With WAVE enabled it also offers WAVE Snow Depth
        Outside Lake, WAVE Water Depth, Water Elevation, and Surface
        Displacement.
  \item \control{WAVE Snow Depth Outside Lake} is a derived copy of AVAC depth
        with AVAC cells intersecting the prepared lake footprint set to zero
        in every band. The original AVAC product is unchanged. Band count and
        elapsed times are preserved.
  \item \control{WAVE Water Elevation} is $B+h$ only on wet cells; dry cells are
        NoData. \control{WAVE Surface Displacement} is $(B+h)$ minus the
        initial surface. Its color ramp is diverging and exactly zero is
        transparent.
  \item \control{Load Time Series} creates or reuses a multiband GeoTIFF and
        loads it into QGIS. AVAC and WAVE layers inherit one shared simulation
        origin, so equal elapsed times align.
  \item The \control{Frame Player} has a seek slider, Play, Pause, Restart, and
        frames-per-second control. It directly switches raster bands and does
        not require tuning QGIS Temporal Controller dates.
  \item \control{Include} selects a simulation-time title, legends, and scale
        bar for PNG output. A legend uses the active layer's QGIS color ramp.
        If multiple visible temporal layers are exported together, each gets
        its own legend.
  \item \control{Extent} chooses Current map canvas or Selected result extent
        for rendering only. \control{Image width} accepts 640--8000 pixels.
  \item \control{Export Current Map PNG} renders the current map and frame.
        \control{Export Time Series Frames} renders the selected series, and
        \control{Export every Nth frame} controls sampling. \control{Cancel
        Export} stops the loop; completed PNG files remain.
\end{itemize}

Temporal and derived rasters are stored under the selected AVAC or WAVE run in
the Working Directory. Frame-export destinations are chosen by the user; the
plugin does not keep generated result rasters in the QGIS application folder.

\uiscreenshot{ui_results_time_series.png}{Time Series Maps with the shared Frame Player and PNG export controls.}

\subsubsection{Profile Analysis}

\begin{itemize}
  \item \control{Profile layer} lists valid QGIS line layers.
        \control{Refresh Layers} refreshes that list.
  \item \control{Selected feature} uses the sole usable line automatically. If
        a layer has multiple lines, select exactly one feature in QGIS; the
        control reports that feature.
  \item \control{Result type} has Selected frame, Full time series (export),
        and Historical maximum through selected frame. Historical maximum is
        calculated independently at every sampled point from frame 1 through
        the current frame; it is not the value of a separate static profile.
  \item \control{Variable} selects any available AVAC, lake-clipped AVAC, or
        WAVE temporal quantity.
  \item \control{Sampling} uses raster-cell spacing or enables a user-defined
        \control{Sample spacing} in meters.
  \item \control{Extract / Plot Profile} samples and plots the chosen result.
        AVAC depth and snow-surface quantities use the same ground/surface
        cross-section presentation as WAVE water profiles. \control{Export
        CSV} writes the current extraction.
  \item For Full time series, \control{Time-series export} writes either one
        long-form CSV or one PNG profile per frame. PNG profile export is
        cancellable.
\end{itemize}

If the selected temporal layer has not been loaded, the profile action loads
it first and then resumes automatically.

\uiscreenshot{ui_results_profile.png}{Profile Analysis. The selected-feature row reflects the actual QGIS line selection.}

\subsubsection{Gauge Histories}

\begin{itemize}
  \item \control{Point layer} accepts any valid QGIS point layer and is empty
        by default. Gauges are postprocessing samples and do not have to be
        defined before AVAC or WAVE preparation.
  \item \control{Variable} independently selects an AVAC or WAVE temporal
        quantity.
  \item \control{Read Gauges from Selected Point Layer} loads a missing result
        series automatically, transforms point coordinates to the result CRS,
        and bilinearly samples every usable point in every band.
  \item \control{Sampled point} lists the sampled features.
        \control{Plot Gauge History} opens the elapsed-time plot and
        \control{Export Gauge CSV} writes simulation time and the selected
        value.
\end{itemize}

\uiscreenshot{ui_results_gauges.png}{Gauge Histories, shared by AVAC and WAVE and based on post-run QGIS point layers.}

\subsubsection{Lake Volume}

This accordion appears only when WAVE is enabled. Lake-water volume is
calculated at every WAVE output time as water depth times cell area, summed
over the initial water-body mask. \control{Plot Lake Volume History} displays
the result and \control{Export Lake Volume CSV} writes elapsed simulation time
and volume in m$^3$. The interface does not currently expose an outflow or
spillway hydrograph.

\uiscreenshot{ui_results_lake_volume.png}{The WAVE-only Lake Volume page in the shared Results tab.}

\clearpage
\section{Working-directory structure}

The Working Directory contains user inputs, isolated runs, derived rasters,
and exports. Managed runtimes are installed separately under the user's
application-support directory.

\begin{verbatim}
AVAC Working Directory/
  derived_inputs/                  point-derived lake polygons
  qgis_previews/                   reusable setup preview rasters
  runs/
    run_YYYYMMDD_HHMMSS/
      .avac_qgis_run.json
      inputs/                      materialized source records and files
      Topo/                        AVAC terrain
      AVAC/
        AVAC_configuration.yaml
        init.avacbin                   packaged binary initial depth
        _output/                   raw AVAC solver output
      qgis_results/                AVAC GeoTIFFs and results.json
  wave_runs/
    wave_YYYYMMDD_HHMMSS/
      .avac_qgis_wave_run.json
      impulse_configuration.yaml
      Topo/                        WAVE terrain and force-dry mask
      CL/                          shoreline faces, source, and source ledger
      Wave/_output/                raw WAVE solver output
      qgis_wave_results/           WAVE GeoTIFFs, volume CSV, manifests
\end{verbatim}

Run directory names and displayed creation times use the computer's local
time. The AVAC and WAVE rasters store elapsed simulation seconds together with
one shared per-simulation temporal origin. A WAVE scenario records a reference
to its completed AVAC source; it does not overwrite, delete, or silently edit
that source run.

\section{Portable installation and documentation source}

The GitHub source archive does not contain compiled managed runtimes and is not
the QGIS installer. Install the platform-specific ZIP from the GitHub Releases
page through Plugins $>$ Manage and Install Plugins $>$ Install from ZIP.

This guide is maintained as
\pathcode{docs/ui\_reference/AVAC\_QGIS\_UI\_REFERENCE.tex}.
Current images are under \pathcode{docs/ui\_reference/images/}; they can be
regenerated from the installed QGIS Python environment with
\pathcode{docs/ui\_reference/generate\_ui\_screenshots.py}. Compile from the
\pathcode{docs/ui\_reference/} directory with \pathcode{latexmk -pdf
AVAC\_QGIS\_UI\_REFERENCE.tex}.

\end{document}
